\chapter*{INTEGRATING MULTIMODAL REPRESENTATIONS INTO GRAPH-BASED FASHION RECOMMENDER SYSTEMS\\ABSTRACT}

Personalized fashion recommender systems are hindered by extreme user-item interaction sparsity and the limitations of ID-bound collaborative filtering. This study integrates visual and textual representations of clothing items into Graph Neural Network (GNN) recommenders. Building on LightGCN as the shared collaborative-filtering backbone, we benchmark three multimodal architectures---CombiGCN (adapted to item similarity graphs), BM3, and FREEDOM---under twenty-four configurations spanning visual encoders (MobileNetV2, CLIP) with BERT-based textual features, and fusion paradigms (late fusion vs. weighted attention gating) on the Vibrent Clothes Rental (VCR) dataset. Three findings emerge. First, within each architecture's best configuration, MobileNetV2 matches or outperforms CLIP, indicating that local texture descriptors suit garment styling; this advantage is configuration-dependent rather than universal, as CLIP surpasses MobileNetV2 in certain visual-only settings. Second, parameter-free late fusion is a stable, robust strategy, whereas trainable attention gating induces overfitting on sparse interaction logs. Third, model superiority is retrieval-depth dependent: BM3, using self-supervised bootstrap alignment, attains the best performance for ranking lists of size $K \ge 5$, while the adapted CombiGCN remains competitive at the shallowest depths. Overall, embedding multimodal features into GNNs helps mitigate data sparsity, offering empirical guidelines for fashion recommendation design.